\documentclass[border=2pt]{standalone}
\usepackage{amsmath, amsthm, amsfonts}
\usepackage{tikz-cd, kotex}
\usepackage{quiver}
\usepackage{Operators}
\usepackage{palette}
\usepackage{diagram-style}
\begin{document}

%%FIG 극한 stable map의 domain과 상 (Intro_Reading_Note-1.svg)
\begin{tikzpicture}
	% domain: main component + bubble
	\draw[thick, black] (0,-2) -- (0,2.3);
	\node at (0,2.6) {main component};
	\draw[thick, black] (0.8,0.6) circle[radius=0.8];
	\node at (0.8,-0.65) {bubble};
	\fill[black] (0,0.6) circle[radius=3pt];
	\node[left] at (-0.15,0.6) {node};
	\fill[accent1] (1.366,1.166) circle[radius=3pt];
	\node[above right, accent1] at (1.4,1.2) {$x_1$};
	\fill[accent2] (1.366,0.034) circle[radius=3pt];
	\node[below right, accent2] at (1.4,0.0) {$x_2$};
	% the map
	\draw[thick, black, -{stealth}] (2.8,0.3) -- (4.4,0.3);
	\node at (3.6,0.7) {$f$};
	% target
	\draw[thick, black] (5.0,-1.7) -- (7.4,1.9);
	\node[right] at (7.4,1.9) {$L$};
	\fill[black] (6.2,0.1) circle[radius=3pt];
	\node[right] at (6.35,0.05) {$p$};
	\node[black!50] at (7.4,-1.6) {$\mathbb{P}^2$};
\end{tikzpicture}

%%FIG 전체 공간의 blow up과 두 section의 분리 (Intro_Reading_Note-2.svg)
\begin{tikzpicture}
	% left: total space P^1 x A^1, coordinates (t,z)
	\draw[thick, black] (0,-2) -- (0,2);
	\node at (0,2.35) {$C_0$};
	\draw[thick, accent1] (-2,0) -- (2,0);
	\node[right, accent1] at (2.05,0) {$z=0$};
	\draw[thick, accent2] (-2,-2) -- (2,2);
	\node[right, accent2] at (2.05,2) {$z=t$};
	\fill[black] (0,0) circle[radius=3pt];
	\node[black!50] at (0,-2.7) {$\mathbb{P}^1\times\mathbb{A}^1$};
	% blow-down arrow
	\draw[thick, black, -{stealth}] (5.6,0) -- (3.6,0);
	\node at (4.6,0.4) {blow up};
	% right: after the blow-up, chart (t,w)
	\begin{scope}[shift={(8.6,0)}]
		\draw[thick, black] (0,-2) -- (0,2);
		\node[right] at (0.15,-1.7) {$E$};
		\draw[thick, black] (0,2) -- (-1.8,2.8);
		\node[above] at (-1.8,2.85) {$\tilde{C}_0$};
		\fill[black] (0,2) circle[radius=3pt];
		\node[right] at (0.15,2.05) {$w=\infty$ (node)};
		\draw[thick, accent1] (-1.8,-0.7) -- (1.8,-0.7);
		\fill[accent1] (0,-0.7) circle[radius=3pt];
		\node[below right, accent1] at (0.1,-0.75) {$w=0$};
		\draw[thick, accent2] (-1.8,0.7) -- (1.8,0.7);
		\fill[accent2] (0,0.7) circle[radius=3pt];
		\node[below right, accent2] at (0.1,0.65) {$w=1$};
		\node[black!50] at (0,-2.7) {$\operatorname{Bl}_{(0,0)}(\mathbb{P}^1\times\mathbb{A}^1)$};
	\end{scope}
\end{tikzpicture}

%%FIG 대각선을 따른 refined Gysin의 cartesian square (Intro_Reading_Note-3.svg)
\begin{tikzcd}
	M_1\times_S M_2 \arrow[r] \arrow[d] & M_1\times M_2 \arrow[d] \\
	S \arrow[r, "\Delta"'] & S\times S
\end{tikzcd}

%%FIG comb 곡선: Y에 잠긴 internal 성분과 지정된 접촉 중복도의 external 성분들 (Intro_Reading_Note-4.svg)
\begin{tikzpicture}
	% Y: the hypersurface, drawn as a line inside ambient X
	\draw[thick, black] (-0.6,0) -- (0.4,0);
	\draw[thick, black] (6.4,0) -- (7.4,0);
	\node[right] at (7.4,0) {$Y$};
	% spine: the internal component, mapped into Y, carrying x_k
	\draw[very thick, accent1] (0.4,0) -- (6.4,0);
	\node[accent1] at (3.4,-1.3) {internal component $\subset Y$};
	% teeth: external components, attached at anonymous nodes (not marked points)
	\draw[thick, black] (1.3,0) .. controls (1.0,0.8) and (1.5,1.4) .. (1.2,2.2);
	\draw[thick, black] (3.5,0) .. controls (3.8,0.8) and (3.3,1.4) .. (3.6,2.2);
	\draw[thick, black] (5.5,0) .. controls (5.2,0.8) and (5.7,1.4) .. (5.4,2.2);
	\node at (3.4,2.6) {external components};
	\fill[black] (1.3,0) circle[radius=3pt];
	\fill[black] (3.5,0) circle[radius=3pt];
	\fill[black] (5.5,0) circle[radius=3pt];
	\node[below] at (1.3,-0.15) {$m^{(1)}$};
	\node[below] at (3.5,-0.15) {$m^{(2)}$};
	\node[below] at (5.5,-0.15) {$m^{(3)}$};
	% x_k: a marked point on the spine, away from the nodes
	\fill[accent2] (2.4,0) circle[radius=3pt];
	\node[below, accent2] at (2.4,-0.15) {$x_k$};
	\node[black!50] at (7.2,2.4) {$X$};
\end{tikzpicture}

\end{document}
